\section[缺陷算子与 Maxwell 理论]{缺陷算子与 Maxwell 理论（4 月 22 日）}

\srcpage{67}

（TODO：今天笔记漏了很多，抱歉。）

今天我们要讲到本课程的主要笑点之一。回忆我们一直在讨论 TFT 中的缺陷算子：$Z(S^{n-1})$ 是局部算子的代数，$Z(S^{n-2})$ 是线算子的幺半范畴，$Z(S^{n-3})$ 是面算子的幺半 2-范畴，等等。当 $n$ 足够大、缺陷的维数足够小时，这些幺半结构是交换的。例如 $n\ge2$ 时局部算子的代数交换。对线算子，交换性有三个层次：幺半（$n=2$）、编织幺半（$n=3$）与对称幺半（$n\ge4$）。这与 Baez--Dolan 的高阶范畴周期表有关。我们特别关心 $n=4$ 中线算子的情形，此时看的是对称幺半范畴 $Z(S^2)$。

局部算子与线算子之间有关系。存在平凡缺陷，它是线算子的对称幺半范畴的单位。于是可以把局部算子放在线内的一个点上，而线上放平凡缺陷。也就是说，局部算子是平凡线算子的自界面。

如果你不喜欢这个版本的图景，可以换个看法：这是球面悬垂（suspension）的图景。例如 $S^{n-2}$ 是一条线的链环，而把 $S^{n-1}$ 切成两个半球，把 $S^{n-1}$ 分解为协边
\srcpage{68}
\[
\varnothing\longrightarrow S^{n-1}\longrightarrow\varnothing,
\]
因此在赤道上放 $Z(S^{n-1})$ 的一个元素就定义了 TODO。跟踪标架或定向很重要，而我们目前不担心这一点；它影响事情的方式是跟踪对偶。（TODO：这里漏了不少，抱歉。）

另一种看法是「意大利饺」或「UFO」图景，它实现等同
\[
D^{n-1}\amalg_{D^{n-1}\setminus\mathrm{pt}}D^{n-1}\cong S^{n-1}.
\]
关键结论是 $Z(S^{n-1})=\End(1_{Z(S^{n-2})})$：局部算子是平凡线算子的自同态。这不只对余维 1 与 2 成立——在更高余维也运行同样的论证。

但回到局部算子与线算子，回忆 $\cM^0_Z:=\Spec Z(S^{n-1})$ 与 $\cM^1_Z:=\Spec Z(S^{n-2})$。把说法略微改写：$\QCoh(\cM^1_Z,\otimes)$ 中的单位是 $\cM^1_Z$ 的结构层。再说一遍，
\[
Z(S^{n-1})=\cO(\cM^0_Z)=\End_{\QCoh(\cM^1_Z)}(\cO)=\Hom(\cO,\cO)=\Gamma(\cO_{\cM^1_Z}).
\]
也就是说，$\cM^0_Z$ 上的函数与 $\cM^1_Z$ 上的整体函数是同一个东西。也就是说，$\cM^0_Z$ 是 $\cM^1_Z$ 的仿射化！（回忆 $X$ 的仿射化是 $\Spec(\cO_X)$。）这是「$Z(S^{n-1})$ 是 $Z(S^{n-2})$ 中单位的自同态」这一事实的代数几何化身。而在拓扑场论中，单位的自同态对应悬垂。

再说一次，这不会止步于余维 2：若 $n$ 足够高，有空间 $\cM^2_Z$ 以及到 $\cM^1_Z$ 的映射。因为我们关心四维，余维 2 最合适：$Z(S^1)$ 不够交换。

（TODO：关于叠的某件事。也许 vacua 模空间要升级为模叠。）

要点：对四维 TFT $Z$，我们造出了 vacua 模空间。用代数几何的话说它该叫「$Z$ 的谱」，但 spectrum 在物理里意味着太多东西，所以更常叫 vacua 模空间。我们本来可以用一堆东西，但我们要用 $\cM^1_Z$。

像上次一样，$Z$ 的谱的一大用处是：用代数几何把 $Z$ 逼近为简单 TFT，特别是上一讲的「B-模型」。你可以问自己：$Z$ 离这些 B-模型有多近？这些 B-模型是 Lagrange 场论，场是到靶 $X$ 的映射，这是谱分解的物理解释。

回忆给定簇或叠 $X$ 与任意维数 $n$，我们造了截断的 $n$ 维 TFT（即赋到维数 $n-1$ 为止的一切）。$B^n_X(M^{n-1})=\cO(\operatorname{Map}(M,X))$——只需要解释 $\operatorname{Map}$ 是什么意思。上次我们说这意味着局部常值映射，因为当 $M$ 是流形而 $X$ 是代数几何对象时，细化这一点很棘手。

簇是从环到集合的某些函子，而叠把集合换成群胚。于是给定流形 $M$ 与叠 $X$，可以更丰富地解释「局部常值」映射：$\operatorname{Map}(\pi_{\le1}M,X)$。这里 $\pi_{\le1}X$ 是 $M$ 的基本群胚。\yizhu{原文此处写作 $\pi_{\le1}X$，但按上下文应为 $\pi_{\le1}M$，即 $M$ 的基本群胚。}{A（原文 p.68）}某种意义上 $M$ 也定义了从环到群胚的函子，但更平凡：把每个环 $R$ 送到 $\pi_{\le1}M$。

我们最喜欢的叠是 $\mathrm{pt}/G$。对这个 $X$，$\operatorname{Map}(M,X)$ 恰好是 $M$ 上 $G$-局部系统的数据！所以这些 B-模型可以比上一讲更有趣。

现在可以把 B-模型推得更远，例如在 $(n-2)$-流形上得到 $\QCoh(\operatorname{Map}(\pi_{\le1}M,X))$，在更高余维流形上得到（高阶）范畴的层，等等。

某种意义上，代数几何内对 $Z$ 的最佳逼近是 $B^n_{\cM^1_Z}$，而 $Z$ 本身是 $\cM^1_Z$-线性的。上次提到，这可以把 $Z$ 变成 $B^{n+1}_{\cM^1_Z}$ 的边界条件。若 $N$ 是 $(n-2)$-流形，$Z(N)$ 是 $Z(S^{n-2})$-模（在一点周围挖去小球）。于是对 $N$ 上所有点得到 $\QCoh(\operatorname{Map}(\pi_{\le1}N,X))$ 在 $Z(N)$ 上的作用，而这个作用是局部常值的。同样，闭 $(n-1)$-流形的态空间上有 $Z(S^{n-1})$ 的局部常值作用，并且它对 $M$ 局部常值。所以我们实际上有 $Z(S^{n-1})^{\#\pi_0(M)}$ 的作用，这意味着不必选基点。也就是说，无需作选择，
\[
\operatorname{Map}(M,\cM^0_Z)=\Spec R^{\otimes\#\pi_0(M)}=(\cM^0_Z)^{\times\pi_0(M)},
\]
所以它实际上是 $\pi_0(M)$ 多个 $\cM^0_Z$ 分支上的层。这是局部算子；对线算子可以做同样的事。这更有趣一点，因为范畴数更高。若有一条从 $x$ 到 $y$ 的路径，就得到这些作用的同构，并得到 $A\cC(\operatorname{Map}(N,\cM^1_Z))$ 在 $Z(S^{n-2})$ 上的作用。\yizhu{原文此处的 $A\cC$ 疑为 $\QCoh$ 之误：上下文一直在用 $\QCoh$ 的作用，且 $A\cC$ 在本讲义中没有定义。}{A（原文 p.68）}继续想这类事情会导向因子化同调。

例如若 $R$ 是交换环、$M$ 是空间，
\[
\int_M R=R^{\otimes\#\pi_0(M)}=\cO(\operatorname{Map}(M,\Spec R)).
\]
这很「无聊」，因为 $R$ 是离散的。若提高范畴数，令 $\cC$ 是幺半范畴，可以得到更有趣的东西：
\[
\int_M \cC=\cO(\operatorname{Map}(M,\Spec\cC)).
\]
在每一点我们都有线算子的作用，而有一种自然的方式把它们组装成相容的对象。

给定 $x\in N$，有赋值映射 $\operatorname{Map}(N,M)\to M$，于是得到拉回映射 $ev_x^*:\QCoh(M)\to\QCoh(\operatorname{Map}(N,M))$。特别地，可以把 $\QCoh(M)$ 的作用拉回为 $\QCoh(\operatorname{Map}(N,M))$ 的作用。而这些作用是局部常值的，所以它们正确拼合。不过要做到这一点，需要知道 $\cC$ 是 Tannakian 的：它是其谱上的层范畴。（TODO：我听对了吗？有点不确定。）

\srcpage{69}

回到拓扑 Maxwell 理论 $\cA_{\mathrm{GL}_1}$。设 $M$ 是闭三维流形，则 $\cA_{\mathrm{GL}_1}(M)$ 是联络叠上局部常值函数的向量空间。若 $\Sigma$ 是闭二维流形，我们赋 $\Pic(\Sigma)$ 上的局部常值层。

于是在每个 $x\in\Sigma$，得到 $\cA_{\mathrm{GL}_1}(S^2)=\mathrm{Loc}(\Pic(\P^1))$ 在 $\cA_{\mathrm{GL}_1}(\Sigma)$ 上的作用。但 $\P^1$ 上的线丛在同构意义下就是 $\Z$——有一些叠结构会在导出层面影响这一点，我们现在不担心。此外张量积对应加法。于是得到 $\mathrm{Loc}(\Z)$，即带张量积的 $\Z$-分次向量空间范畴。这是 Gauss 定律的一个版本：更大球内有两个小球，它们的「通量」（flux）相加，通量是这些线丛的 Chern 类。而 Chern 类相加。

我们看到 $\cM^0_\cA$ 平凡，但 $\cM^1_\cA$ 是 $\mathrm{pt}/\G_m$，即线丛的模叠。一般地 $\operatorname{Map}(N,\cM^1_\cA)=\mathrm{Loc}_1(N)$。于是发现 $\cA(\Sigma)$ 被 $\QCoh(\mathrm{Loc}_1(\Sigma))$ 作用，即被 $\QCoh(B(\Sigma))$ 作用。于是在对电磁对偶一无所知的情况下，我们从 A-模型里跳出了 B-模型。如果有足够的场论，就不必想 Fourier 变换。

现在你说，「被作用」不是我们承诺的东西，我们承诺的是等价。但结果是一个是另一个上的秩一自由模，所以只需指定一个生成元。

下次我们会在非交换情形讨论这件事，从 $\cA_G(\Sigma)=\mathrm{Sh}(\Bun_G(\Sigma))$ 出发并构造谱侧；几何 Langlands 猜想就是这两侧等价。
